% "The 20-Watt Computer": the popular ("take-along") version.
% Each chapter paperback/NN.tex corresponds to the full chapter ../chapters/NN_*.tex with the same number;
% on the website the reader switches a chapter between the "Short" and "Full" versions.
\documentclass[11pt,openany]{book}
\usepackage[paperwidth=13cm,paperheight=20cm,top=1.8cm,bottom=1.8cm,inner=1.6cm,outer=1.4cm]{geometry}
\usepackage[T2A]{fontenc}
\usepackage[utf8]{inputenc}
\usepackage[english,ukrainian]{babel}
\usepackage{amsmath,amssymb}
\usepackage{xcolor}
\usepackage[unicode,colorlinks=true,linkcolor=blue!60!black]{hyperref}
\usepackage{microtype}
\input{paperback/pencil} % minimal colour-pencil illustrations, one per chapter
\setlength{\parskip}{0.3em}
\setlength{\emergencystretch}{2em} % narrow paperback page: allow a little extra stretch
\renewcommand{\chaptermark}[1]{\markboth{\small #1}{}}
% neuron type code: first four letters of the primary mediator (as in the full edition)
\newcommand{\nt}[1]{\textsf{\textbf{#1}}}
% pointer to the full edition at the end of each chapter
\newcommand{\fullversion}[1]{\par\bigskip\noindent\small\emph{Формули, моделі, задачі та джерела --- у повній версії, розділ~#1.}\normalsize}

\title{\Huge Комп'ютер на 20 ватах\\[14pt] \Large Як обчислює мозок --- коротко і просто}
\hypersetup{pdftitle={Комп'ютер на 20 ватах. Як обчислює мозок --- коротко і просто},pdfauthor={Костянтин Кладько}}
\author{\Large Костянтин~Кладько}
\date{}

\begin{document}
\frontmatter
\input{paperback/cover} % cover: a collage of the chapter drawings
% copyright page (same licence as the full edition)
\thispagestyle{empty}
\vspace*{\fill}
\noindent{\footnotesize \copyright~2026 Костянтин~Кладько.\\[4pt]
Текст і рисунки поширюються за ліцензією Creative Commons «Зазначення авторства --- Некомерційна --- Поширення на тих самих умовах» 4.0 (CC BY-NC-SA 4.0): \url{https://creativecommons.org/licenses/by-nc-sa/4.0/deed.uk}. Книгу можна вільно копіювати, використовувати у викладанні, перекладати й доопрацьовувати з некомерційною метою за умови зазначення автора та поширення похідних творів на тих самих умовах.\\[4pt]
Повна версія, вихідні тексти та моделі: \url{https://github.com/kladkogex/20-watt-computer}.}
\clearpage
\tableofcontents
\input{paperback/00_preface}
\mainmatter
\input{paperback/01}
\input{paperback/02}
\input{paperback/03}
\input{paperback/04}
\input{paperback/05}
\input{paperback/06}
\input{paperback/07}
\input{paperback/08}
\input{paperback/09}
\input{paperback/10}
\input{paperback/11}
\input{paperback/12}
\input{paperback/13}
\input{paperback/14}
\input{paperback/15}
\input{paperback/16}
\input{paperback/17}
\input{paperback/18}
\input{paperback/19}
\input{paperback/20}
\input{paperback/21}
\input{paperback/22}
\end{document}
